% !PRINTED 503
\weberpage{503}

To establish the theorem completely, we must still prove that the symmetric functions of $M_1$, $M_2$, $M_3$ are, in
the exceptional case last considered, rational functions of
$j(\omega)$. This requires us to investigate what becomes of (6)
when, for a singular value $u$ of $v$, the denominator $F'(v)$ vanishes, and,
since $M$ remains finite, the numerator as well. We again let
$\omega$ be variable to begin with, and denote the roots of
(4) by

\weq{11}{v_1,\, v_2,\, v_2,\, \ldots,\, v_\nu.}
% [sic] printed repeats "v_2" in the list — presumably an intended index run
% $v_1,\ldots,v_\nu$
Now let, for the particular value of $\omega$ under consideration,

\weq{12}{v_1,\, v_2,\, \ldots,\, v_\lambda}
become equal to one another, while the remaining $v_{\lambda+1}, \ldots, v_\nu$ stay
different from them. We take any symmetric function of the quantity (12), for example, for an indeterminate $t$:

\weq{13}{\sigma = (t - v_1),\, (t - v_2),\, \ldots,\, (t - v_\lambda);}
% [sic] printed separates the factors with commas rather than a product sign
if we carry out in $\sigma$ all permutations of the quantities (11), we
determine $r$ values

\weq{14}{\sigma_1,\, \sigma_2,\, \ldots,\, \sigma_r,}
when
\[
r = \frac{\nu!}{\lambda!\,(\nu - \lambda)!}
\]
denotes the number of combinations of $\nu$ quantities taken $\lambda$ at a time (without repetition). We can dispose of $t$ in such a way that, even for
the singular value $\omega$, \gesperrt{only one} of the quantities (14) becomes equal to the
first $\sigma$. The $\sigma$ are the roots of an equation
\[
\Phi(x, u) = 0
\]
of degree $r$, and for the singular value $x = \sigma$ the quantity $\Phi'(\sigma)$
remains different from zero.

Let us now consider a symmetric function $S$ of the quantity
corresponding to (12)
\[
M_1,\, M_2,\, \ldots,\, M_\lambda,
\]
then, through the $\nu$ combinations of the numbers $a$, $\partial$, $c$ that
lead to the values (11), it likewise takes $r$ values $S_1, S_2, \ldots, S_r$,
and the sum
\[
\Phi(x, u) \sum \frac{S}{x - \sigma} = \Theta(x, u)
\]
is a rational integral function of $x$ and $u$. From this it follows,
on letting $x$ pass into $\sigma$:
\[
S = \frac{\Theta(\sigma, u)}{\Phi'(u)}.
\]
% [sic] German source prints denominator $\Phi'(u)$; the preceding specialization $x=\sigma$ requires $\Phi'(\sigma)$
